\section{An exact finite experiment}\label{ch14:sec:experiment}
A sign sequence of length $n$ is fixed by $n$ independent choices, one sign for each position, so there are $2^n$ sign sequences of length $n$. For small $n$ a computer can examine every one of them. The table below shows the result of such an exhaustive search for $1\le n\le13$. Only integers are involved, so the arithmetic is exact, with no rounding. Appendix~\ref{app:practical} contains a short program that produces the table directly from the definitions. For each length the search counts two kinds of sequences:
\begin{itemize}
\item the Barker sequences;
\item the sequences with $C_s=0$ for every shift $s$ with $1\le s\le n-1$, which by Section~\ref{ch14:sec:circulant} are exactly the first rows of circulant Hadamard matrices.
\end{itemize}
Every list is counted on its own. For instance, a sequence and its negative are counted as two sequences. So are a sequence and its reversal (the same entries in the opposite order), and a sequence and its rotations, whenever these lists are different.

\begin{center}
\begin{tabular}{@{}r rr@{}}\toprule
Length $n$ & Barker sequences & Sequences with $C_s=0$ for $1\le s\le n-1$\\\midrule
1&2&2\\2&4&0\\3&4&0\\4&8&8\\5&4&0\\6&0&0\\7&4&0\\8&0&0\\9&0&0\\10&0&0\\11&4&0\\12&0&0\\13&4&0\\\bottomrule
\end{tabular}
\end{center}

The table is evidence about these thirteen lengths and no others. It cannot rule out a qualifying sequence of some larger length that was never tested. Nor can the search simply be pushed much further, because each additional position doubles the number of sequences. Length $60$ alone has $2^{60}$ sign sequences, which is more than $10^{18}$. A computer that examined a billion sequences per second would need more than $36$ years for that single length. More computing power is therefore not a way to settle the question for every length. Only a proof can do that.

Every pattern in the table agrees with what we have proved, and each one has a reason.
\begin{itemize}
\item The last column is zero at every length except $1$ and $4$. Proposition~\ref{thm:circulantnecessary} predicts this: below $16$, the only orders it allows are $1$ and $4$, and at both of these we already know examples.
\item For the odd lengths greater than $1$ there is also a simpler explanation, given in the pause at the end of Section~\ref{ch14:sec:periodic}: when $n$ is odd, every periodic correlation is odd, so none of them can be zero.
\item Every count in the table is even. By Exercise~\ref{ex:14.2}, negating every sign changes none of the correlations, aperiodic or periodic, because each correlation is a sum of pair products and pair products do not change (Section~\ref{ch06:sec:wlog}). So whenever a sequence $a$ qualifies, its negative $-a$ qualifies too. Moreover $-a$ is never equal to $a$, since that would require $a_0=-a_0$, that is, $a_0=0$. So the qualifying sequences of each length split into pairs $\{a,-a\}$, and their number is even.
\end{itemize}
The normalization of Section~\ref{ch06:sec:wlog} would allow a faster search: examine only the sequences that begin with $1$, and double each count. The program in Appendix~\ref{app:practical} examines all $2^n$ sequences instead, so the evenness of every count is a small independent check on its output.
